\subsubsection{Tempo Médio Real de Sessão}

O tempo médio real de sessão por aluno é calculado com base nos \textit{traces} de interação coletados pela extensão durante a aula. O algoritmo implementado no Script segue as seguintes etapas:

\begin{enumerate}
    \item \textbf{Agrupamento por aluno}: Todos os \textit{traces} são organizados por identificador único de usuário, criando uma lista temporal de interações para cada participante.
    
    \item \textbf{Processamento individual}: Para cada aluno, os \textit{timestamps} são convertidos para objetos de tempo e ordenados cronologicamente.
    
    \item \textbf{Cálculo de duração}: A diferença entre o primeiro e último registro de interação representa o tempo total de sessão do aluno, convertido para minutos.
    
    \item \textbf{Valores mínimos}: Casos com apenas um registro recebem duração mínima de 1,0 minuto, enquanto sessões calculadas com menos de 1 minuto são ajustadas para este piso.
    
    \item \textbf{Média agregada}: A métrica final é obtida através da média aritmética simples de todas as durações individuais de sessão.
\end{enumerate}

% \lstinputlisting[
%     language=python,
%     caption={Cálculo do Tempo Médio Real de Sessão},
%     firstline=4,
%     lastline=62,
%     label={lst:calculate_real_avg_session_per_student}
% ]{scripts/metricsCalc/base.py}



\textbf{Premissas do cálculo}:
\begin{itemize}
    \item A extensão permanece ativa continuamente durante a participação do aluno
    \item Os \textit{traces} são emitidos regularmente durante a sessão
    \item O primeiro registro representa o início efetivo da participação
    \item O último registro representa o término da sessão do aluno
\end{itemize}

Esta métrica fornece uma estimativa do tempo real que cada aluno dedicou à aula, considerando sua presença em sala de aula virtual.

\subsubsection{Taxa de Presença}
A Taxa de Presença (\textit{Attendance Ratio}) representa a proporção entre o tempo médio real de participação dos alunos e a duração total oficial da aula. Esta métrica quantifica o nível de engajamento temporal da turma como um todo, indicando qual porcentagem da aula os alunos permaneceram ativos em média.

O cálculo é realizado pela função \textbf{calculate\_attendance\_ratio} (Script~\ref{lst:calculate_attendance_ratio}) conforme a seguinte fórmula:

\begin{equation}
\text{Taxa de Presença} = \min\left(1, \max\left(0, \frac{\text{Tempo Médio Real de Sessão}}{\text{Duração Oficial da Aula}}\right)\right)
\end{equation}


% \lstinputlisting[
%     language=python,
%     caption={Cálculo da Taxa de Presença},
%     firstline=104,
%     lastline=129,
%     label={lst:calculate_attendance_ratio}
% ]{scripts/metricsCalc/base.py}


\textbf{Etapas do cálculo}:
\begin{enumerate}
    \item \textbf{Entrada}: Recebe os \textit{traces} de interação e o identificador da aula
    \item \textbf{Tempo médio real}: Calcula o tempo médio de sessão usando a função anterior
    \item \textbf{Duração oficial}: Obtém o tempo total da aula do sistema de registros
    \item \textbf{Razão}: Divide o tempo médio real pela duração oficial
    \item \textbf{Normalização}: Limita o resultado entre 0 (0\%) e 1 (100\%)
    \item \textbf{Arredondamento}: Retorna o valor com 3 casas decimais de precisão
\end{enumerate}

